\section{Background}
\label{sec:background}

This section introduces the EML function and its universality properties, describes the
verified primitive constructions that form the basis of our method, and reviews the two
families of prior work---physics-informed neural networks and symbolic
regression---against which we position our approach.

%% ---------- 2.1 ----------
\subsection{The EML Function}
\label{sec:eml}

The \emph{Exp-Minus-Log} (EML) function is defined as
\begin{equation}
  \mathrm{eml}(x,y) \;=\; e^{x} - \ln y, \qquad x \in \mathbb{R},\; y > 0.
  \label{eq:eml-def}
\end{equation}
Odrzywo{\l}ek~\cite{odrzywołek2026} established that $\mathrm{eml}$ is a
\emph{universal function}: every computable function admits a representation as a
finite recursive composition of $\mathrm{eml}$ operations.  This result is
surprising in that a single bivariate primitive suffices to generate the full
hierarchy of elementary and special functions.

Several elementary identities follow directly from the definition:
\begin{align}
  \mathrm{eml}(x, 1) &= e^{x},
    &\text{since } \ln 1 = 0, \label{eq:eml-exp}\\
  \mathrm{eml}(0, y)  &= 1 - \ln y,
    &\text{since } e^{0} = 1, \label{eq:eml-const}\\
  \frac{\partial\,\mathrm{eml}}{\partial x} &= e^{x},
    &\frac{\partial\,\mathrm{eml}}{\partial y} = -\frac{1}{y}.
    \label{eq:eml-deriv}
\end{align}
The closure of $\mathrm{eml}$ under differentiation---each partial derivative
is itself expressible via the primitives $e^{x}$ and $1/y$---is a property we
exploit heavily in Section~\ref{sec:method}, where PDE residuals must be
differentiated through the compositional graph.

\paragraph{Limitations.}
The universality claim is subject to important caveats.
Smith~\cite{smith2024} demonstrated that $\mathrm{eml}$ cannot express roots of
generic quintic polynomials whose associated Galois groups are non-solvable,
a consequence of the Abel--Ruffini theorem applied to the monodromy structure of
$\mathrm{eml}$ compositions.  Separately, \.{I}pek~\cite{ipek2024}
proposed an extension of the universality result (arXiv:2604.13871), but the paper
was subsequently withdrawn by the author citing a ``fundamental limitation'' in the
proof.  These results constrain the theoretical scope of EML universality but do
not affect the practical utility of the framework for compositions of elementary
functions---the regime relevant to PDE solutions encountered in physics.

\paragraph{Depth requirements.}
Even for elementary operations, the depth of the required $\mathrm{eml}$ composition
can be substantial.  Odrzywo{\l}ek~\cite{odrzywołek2026} reports that expressing
$\ln$ requires depth~3, multiplication requires a reverse-Polish notation (RPN)
representation of length~41, and the constant $\pi$ requires length~193 or more.
These figures render direct gradient-based optimization over raw EML expression
trees impractical: the search space is combinatorially vast, and the resulting
computation graphs are deep enough to suffer from vanishing and exploding
gradients.  This observation motivates the strategy described in
Section~\ref{sec:primitives}: rather than learning $\mathrm{eml}$ trees, we
pre-compute algebraically verified constructions for a small set of elementary
functions and optimize over compositions of those constructions.

%% ---------- 2.2 ----------
\subsection{Verified Primitive Constructions}
\label{sec:primitives}

The central idea of our method is to bridge the gap between the theoretical
universality of $\mathrm{eml}$ and the practical requirements of gradient-based
optimization.  We achieve this by pre-computing a library of fixed,
algebraically verified constructions that express standard elementary functions
as $\mathrm{eml}$ compositions.  Each construction is implemented as a frozen
PyTorch module with no learnable parameters; the optimization problem is thus
lifted from the space of raw $\mathrm{eml}$ trees to the space of linear
combinations of these verified primitives.

\paragraph{Exponential.}
The simplest construction is immediate from~\eqref{eq:eml-exp}:
\begin{equation}
  \exp(x) = \mathrm{eml}(x, 1).
  \label{eq:prim-exp}
\end{equation}
This is a depth-1 composition.

\paragraph{Natural logarithm.}
The logarithm requires depth~3.  We construct it as follows:
\begin{equation}
  \ln(z) = \mathrm{eml}\!\Big(1,\;\mathrm{eml}\!\big(\mathrm{eml}(1,z),\;1\big)\Big).
  \label{eq:prim-ln}
\end{equation}
To verify, evaluate from the inside outward.  First,
$\mathrm{eml}(1, z) = e - \ln z$.  Second,
$\mathrm{eml}(e - \ln z,\; 1) = \exp(e - \ln z)$, since $\ln 1 = 0$.
Third,
\[
  \mathrm{eml}\!\big(1,\; \exp(e - \ln z)\big)
  = e - \ln\!\big(\exp(e - \ln z)\big)
  = e - (e - \ln z)
  = \ln z.
\]

\paragraph{Sine and cosine.}
The trigonometric functions are obtained via the complex extension of
$\mathrm{eml}$.  Because $\mathrm{eml}(ix, 1) = \exp(ix) - \ln 1 = \exp(ix)$,
Euler's formula gives
\begin{align}
  \sin(x) &= \operatorname{Im}\!\big(\mathrm{eml}(ix, 1)\big),
    \label{eq:prim-sin}\\
  \cos(x) &= \operatorname{Re}\!\big(\mathrm{eml}(ix, 1)\big).
    \label{eq:prim-cos}
\end{align}
In practice, we implement $\mathrm{eml}(ix, 1)$ using PyTorch's complex tensor
arithmetic, extracting the real and imaginary parts to obtain $\cos$ and $\sin$
respectively.

\paragraph{The constant $\boldsymbol{\pi}$.}
The circle constant is recovered from the complex logarithm construction:
\begin{equation}
  \pi = -\operatorname{Im}\!\big(\ln(-1)\big),
  \label{eq:prim-pi}
\end{equation}
where $\ln$ is the depth-3 EML construction of~\eqref{eq:prim-ln} extended to
complex arguments.

\paragraph{Verification protocol.}
It is essential to emphasize that none of these constructions are learned.
They are algebraic identities, proved symbolically and implemented as
deterministic modules.  We verify each implementation numerically against the
corresponding PyTorch intrinsic (\texttt{torch.sin}, \texttt{torch.cos},
\texttt{torch.exp}, \texttt{torch.log}) over $10^{6}$ uniformly sampled
points, confirming pointwise agreement to within $10^{-5}$.  This verification
step guarantees that the compositional basis inherits the full numerical
fidelity of the underlying hardware transcendentals.

%% ---------- 2.3 ----------
\subsection{Physics-Informed Neural Networks}
\label{sec:pinns}

Physics-informed neural networks (PINNs), introduced by Raissi, Perdikaris,
and Karniadakis~\cite{raissi2019}, embed the governing PDE directly into the
training loss.  Given a PDE of the form $\mathcal{N}[u] = 0$ with boundary
conditions $\mathcal{B}[u] = 0$, a PINN parameterizes the solution as a neural
network $u_{\theta}$ and minimizes
\begin{equation}
  \mathcal{L}
  = \underbrace{\frac{1}{N_d}\sum_{i=1}^{N_d}
      \big|u_{\theta}(\mathbf{x}_i) - u_i^{\mathrm{data}}\big|^{2}
    }_{\mathcal{L}_{\mathrm{data}}}
  + \lambda\,
    \underbrace{\frac{1}{N_r}\sum_{j=1}^{N_r}
      \big|\mathcal{N}[u_{\theta}](\mathbf{x}_j)\big|^{2}
    }_{\mathcal{L}_{\mathrm{PDE}}},
  \label{eq:pinn-loss}
\end{equation}
where $\{\mathbf{x}_i, u_i^{\mathrm{data}}\}$ are observed data points,
$\{\mathbf{x}_j\}$ are collocation points sampled in the interior of the
domain, and $\lambda$ is a weighting hyperparameter.

PINNs possess several attractive properties: they are mesh-free, accommodate
irregular geometries naturally, and can incorporate partial or noisy
observations.  These features have driven widespread adoption across
computational physics.  However, for the purposes of the present work, PINNs
suffer from a fundamental limitation: their output is a set of network weights
$\theta$, not a symbolic expression.  Consequently, PINN solutions offer no
direct interpretability, cannot be analytically verified against the governing
equations, and do not transfer to related problems in the way that a closed-form
expression does.  Our method retains the physics-informed loss
structure~\eqref{eq:pinn-loss} but replaces the neural network ansatz with a
compositional search over EML-derived primitives, thereby recovering the
symbolic form of the solution.

%% ---------- 2.4 ----------
\subsection{Symbolic Regression}
\label{sec:symreg}

Symbolic regression seeks to identify a mathematical expression that fits
observed data, searching simultaneously over functional forms and numerical
coefficients.  Two recent systems are particularly relevant.

\textbf{PySR}~\cite{cranmer2023} combines genetic programming with local
gradient-based refinement.  An evolving population of expression trees is
subject to mutation, crossover, and selection operators, with periodic
gradient descent applied to the numerical constants in each tree.  PySR has
demonstrated strong performance on a range of scientific benchmarks.

\textbf{AI Feynman}~\cite{udrescu2020} takes a complementary approach,
exploiting physics-inspired strategies such as dimensional analysis, symmetry
detection, and separability testing to decompose complex expressions into
simpler sub-problems before applying brute-force or neural-network-guided
search.

Both methods excel at recovering algebraic laws of moderate complexity---the
paradigmatic examples being $F = ma$ and $E = mc^{2}$.  However, they face
significant challenges in the regime targeted by the present work:
\begin{itemize}
  \item \emph{Multi-variable PDE solutions.}  Solutions to PDEs in two or more
    spatial dimensions typically involve products of univariate functions along
    different axes (e.g., $\sin(x)\cos(y)$).  Standard symbolic regression
    treats the entire expression as a single tree, making it difficult to
    exploit this separable structure.
  \item \emph{Combinatorial explosion.}  The space of nested function
    compositions grows super-exponentially with expression depth.  Genetic
    operators struggle to navigate this space efficiently when the target
    involves three or more levels of composition.
  \item \emph{Mixed exponential--trigonometric forms.}  Time-dependent
    solutions with spatial oscillation and temporal decay---such as
    $\sin(x)\exp(-\alpha t)$---require the simultaneous discovery of
    exponential and trigonometric primitives with coupled coefficients, a
    combination that lies in a particularly sparse region of typical
    expression-tree search spaces.
\end{itemize}

Our approach addresses these difficulties by constraining the search to a
physically motivated basis of EML-derived primitives organized into separable,
axis-aligned terms.  Rather than evolving arbitrary expression trees, we
optimize a structured ansatz whose functional building blocks---$\sin$, $\cos$,
$\exp$, $\ln$---are fixed and verified, while only the linear combination
coefficients and frequency parameters remain trainable.  This reduces the
effective search space by orders of magnitude while retaining the capacity to
express the product-of-functions structure characteristic of classical PDE
solutions.
